\documentclass[10pt,a4paper]{amsart}
\usepackage[margin=19mm]{geometry}
\usepackage{amsmath,amssymb,mathtools}
\usepackage[scr=rsfs,cal=euler]{mathalfa}
\usepackage{xfrac}
\usepackage[unicode,hidelinks,bookmarks=false]{hyperref}
\newcommand{\ie}{i.e.\ }
\newcommand{\cf}{cf.\ }
\newcommand{\resp}{respectively\ }
\newcommand{\Subsection}{\subsection}
\providecommand{\nobreakdash}{\nobreak\hbox{-}}
\setlength{\emergencystretch}{2em}
\multlinegap0pt
\usepackage{amssymb}
\usepackage{mathtools}
\usepackage{algorithm}\usepackage{algpseudocode}
\usepackage{xcolor}
\newtheorem{lemma}{Lemma}
\newcommand*{\ZZ}{\mathbb Z}
\title{The threshold for online balancing of i.i.d. binary vectors}
\author{Dylan J. Altschuler, Konstantin Tikhomirov}
\date{Source: arXiv:2609.14975v1 (CC BY 4.0)}
\begin{document}
\maketitle

\noindent\emph{Source and licence.} The threshold for online balancing of i.i.d. binary vectors. Version arXiv:2609.14975v1 (CC BY 4.0), distributed by arXiv under the Creative Commons Attribution 4.0 International licence, http://creativecommons.org/licenses/by/4.0/, as recorded in the arXiv licence metadata for this exact version and shown on the abstract page https://arxiv.org/abs/2609.14975v1. Full licence notice: \texttt{licenses/arxiv-2609.14975-CC-BY-4.0.txt}.\par\smallskip

\noindent\textit{Editorial scope.} Counted unit: the article's lemma lem:potential (source0.tex lines 957--985). The article's complete original proof of it is delivered with it (source0.tex lines 987--1104, one \texttt{proof} environment of 118 lines). No mathematics is written, repaired or re-derived: the statement and the proof are the article's own lines, in the article's own notation, and no retained line is substituted or re-flowed.  No further source-local context is needed: every cross-reference of every retained line resolves inside the two delivered spans.  Nothing was omitted from any retained span, so the package carries no omission record.  No retained line contains a citation, so the wrapper carries no bibliography and no bibliography entry.  No retained line contains a figure environment, an image-inclusion command or drawing code, so no image is delivered and the package declares no asset dependency.  Editorial formatting only, in addition: an AMS \texttt{amsart} preamble that restates the article's own declarations and macros used by the retained lines, namely the environments \texttt{lemma} on separate counters, together with the standard packages those lines need.  The retained environments print in this document as Lemma 1 in order of appearance, whereas the article prints the same items with the numbering of its own full document.  The complete native source of the article as distributed in the arXiv e-print for this exact version is bundled unchanged as \texttt{sources/210380/source0.tex}.
\begin{lemma}[Finite-difference estimates]\label{lem:potential}
Assuming $a$ is sufficiently large, the following estimates hold with
absolute constants $c,C>0$:
\begin{equation}\label{eq:lambda-theta-bounds}
    \frac{c}{a}\le\lambda\le\frac{C}{a},
    \qquad
    \Theta\le Ca(\sqrt d+\log Y).
\end{equation}
For every $k\in\ZZ$,
\begin{equation}\label{eq:first-difference-bound}
    |\nabla F(k)|\le\frac{C}{a}F(k).
\end{equation}
Moreover,
\begin{align}
 (\Delta F(k))_+&\le\frac{C}{a^2d},
       && |k|\le\rho/2,\label{eq:second-difference-small}\\
 (\Delta F(k))_+&\le\frac{C}{a}|\nabla F(k)|,
       && |k|>\rho/2,\label{eq:second-difference-large-1}\\
 (\Delta F(k))_+&\le
       C|\nabla F(k)|\sqrt{\frac{F(k)}{a^2d}},
       && |k|>\rho/2.\label{eq:second-difference-large-2}
\end{align}
Finally,
\begin{equation}\label{eq:second-by-first}
 |\Delta F(k)|\le|\nabla F(k)|\quad(k\ne0),
 \qquad
 |\Delta F(0)|\le\frac{C}{a^2d}.
\end{equation}
\end{lemma}

\begin{proof}
By evenness, it suffices to consider nonnegative integers $k$.
Throughout the proof, $c,C>0$ are absolute constants, and we take $a\ge4$.
Let
\[
    g(x):=\left(\frac{\rho^2}{\rho^2-x^2}\right)^2,
    \qquad |x|<\rho.
\]
Since $a\le\rho-b<a+1$, we have
\[
 \frac{1}{2a}
 \le\frac{2b+1}{(\rho-b-1)(\rho+b+1)}
 \le\frac{3}{a},
 \qquad
 \frac{d}{16}\le g(b)\le g(b+1)\le16d.
\]
The identity
\[
 \lambda=2\log\left(1+
   \frac{2b+1}{(\rho-b-1)(\rho+b+1)}\right)
\]
therefore gives $1/(2a)\le\lambda\le6/a$.
In the range $2\le d\le {n^{1/5}}$, we have $Y>16d$, so
$\Theta\ge b+2$ and
\[
 \Theta=b+\left\lceil\lambda^{-1}\log\frac{Y}{g(b)}\right\rceil
 \le Ca(\sqrt d+\log Y).
\]
This proves \eqref{eq:lambda-theta-bounds}. We now treat the inverse-square,
{\color{black}hyperbolic-cosine}, and truncated parts of the potential in turn.

For $0\le k\le b$, the three values $F(k-1),F(k),F(k+1)$ agree with
$g(k-1),g(k),g(k+1)$, including at $k=b$ because
$g(b)e^\lambda=g(b+1)$. Differentiating gives
\begin{equation}\label{eq:g-derivatives}
 \frac{g'(x)}{g(x)}=\frac{4x}{\rho^2-x^2},
 \qquad
 \frac{g''(x)}{g(x)}
   =\frac{4\rho^2+20x^2}{(\rho^2-x^2)^2}.
\end{equation}
Since $\rho-k\ge a\ge4$, these formulas imply, for $|x-k|\le1$,
\begin{equation}\label{eq:g-local-bounds}
 \frac{g(k)}{C}\le g(x)\le Cg(k),\qquad
 |g'(x)|\le\frac{Cg(k)}{\rho-k},\qquad
 0\le g''(x)\le\frac{Cg(k)}{(\rho-k)^2}.
\end{equation}
Using the integral identities
\[
 \nabla F(k)=\frac12\int_{-1}^{1}g'(k+t)\,dt,
 \qquad
 \Delta F(k)=\frac12\int_{-1}^{1}(1-|t|)g''(k+t)\,dt,
\]
we obtain
\[
 |\nabla F(k)|\le\frac{C}{a}F(k),
 \qquad
 0\le\Delta F(k)\le\frac{CF(k)}{(\rho-k)^2}.
\]
If $k\le\rho/2$, then $F(k)\le16/9$ and $\rho-k\ge\rho/2$, proving
\eqref{eq:second-difference-small}. If $\rho/2<k\le b$, then
\eqref{eq:g-derivatives} also gives
$g'(x)\ge cg(k)/(\rho-k)$ for $|x-k|\le1$. Consequently,
\[
 \Delta F(k)\le\frac{C}{\rho-k}|\nabla F(k)|.
\]
Both remaining second-difference bounds follow from
\[
 \frac{1}{\rho-k}\le\frac1a,
 \qquad
 \sqrt{F(k)}=\frac{\rho^2}{(\rho-k)(\rho+k)}
 \ge\frac{a\sqrt d}{2(\rho-k)}.
\]

For $b<k\le\Theta-2$, {\color{black}the three potential values $F(k-1),F(k),F(k+1)$ satisfy $F(k\pm1)=e^{\pm\lambda}F(k)$.} Hence
\[
 \nabla F(k)=F(k)\sinh\lambda,
 \qquad
 \Delta F(k)=F(k)(\cosh\lambda-1).
\]
Since $\lambda\le6/a$, it follows that
\[
 |\nabla F(k)|\le\frac{C}{a}F(k),
 \qquad
 \Delta F(k)=\tanh(\lambda/2)|\nabla F(k)|
 \le\frac{C}{a}|\nabla F(k)|.
\]
Moreover, $F(k)\ge g(b)\ge d/16$, so
$1/a\le4\sqrt{F(k)/(a^2d)}$. This also proves
\eqref{eq:second-difference-large-2} in this range.


At $k=\Theta-1$, truncation only decreases the forward increment
$u:=F(k+1)-F(k)$, leaving $v:=F(k)-F(k-1)>0$ unchanged.
Both $\nabla F(k)=(u+v)/2$ and
\[
 \frac{(\Delta F(k))_+}{\nabla F(k)}
 =\frac{(u-v)_+}{u+v}
\]
are nondecreasing in $u$. Thus the bounds just proved for the {\color{black}hyperbolic-cosine} potential remain valid here. At $k=\Theta$,
\[
 \Delta F(k)=-\nabla F(k),\qquad
 0\le\nabla F(k)=\frac{Y-F(k-1)}2
 \le\frac{e^\lambda-1}{2}F(k-1)\le\frac{C}{a}F(k).
\]
All differences vanish for $k>\Theta$. This completes the proof of
\eqref{eq:first-difference-bound}--\eqref{eq:second-difference-large-2}.


Finally, for $k\ge1$, let
$u:=F(k+1)-F(k)\ge0$ and $v:=F(k)-F(k-1)\ge0$.
Then
\[
 |\Delta F(k)|=\frac{|u-v|}{2}
 \le\frac{u+v}{2}=|\nabla F(k)|.
\]
At $k=0$, the bound follows from \eqref{eq:second-difference-small}.
This proves \eqref{eq:second-by-first}.
\end{proof}

\end{document}
